\documentclass[11pt]{article}
\usepackage[a4paper,margin=2.4cm]{geometry}
\usepackage{amsmath,amssymb,booktabs,microtype}
\usepackage[colorlinks=true,linkcolor=black,urlcolor=blue]{hyperref}

\newcommand{\Vset}{\mathcal{V}}
\newcommand{\Tset}{\mathcal{T}}
\newcommand{\relu}[1]{\left(#1\right)^{+}}

\title{\textbf{DEX-AR: the equations, before and after the rank prefilter}}
\author{}
\date{}

\begin{document}
\maketitle
\vspace{-2.2\baselineskip}

\noindent\rule{\textwidth}{0.4pt}
\vspace{0.4em}

\noindent This note states what DEX-AR computes as published, then states the one
line that changes. Everything here is transcribed from the implementation in
\texttt{dexar/wrapper.py}; equation numbers in parentheses refer to the paper
(Bousselham et al., \emph{DEX-AR: A Dynamic Explainability Method for
Autoregressive VLMs}, arXiv:2603.06302).

\section{Notation}

For a generation step $t$, layer $l \in \{1,\dots,L\}$ and attention head
$i \in \{1,\dots,H\}$:

\begin{tabular}{@{}ll@{}}
\toprule
$\Vset$, $\Tset$ & the image-token and text-token positions of the sequence; $N = |\Vset| = h \times w$ \\
$A^{l,t}$        & attention probabilities at layer $l$, step $t$ \\
$y^{l}_{t}$      & hidden state at the last position, layer $l$ \\
$o^{l,t}$        & intermediate logits, $o^{l,t} = \mathrm{LM\_Head}(y^{l}_{t})$ \quad (Eq.~3, the logit lens) \\
$v_t$            & the sampled token at step $t$: the final layer's argmax \quad (Eq.~2) \\
\bottomrule
\end{tabular}

\medskip
\noindent DEX-AR differentiates the sampled token's intermediate logit with
respect to the attention probabilities, and keeps only the last query row --- the
row belonging to the position that is about to emit $v_t$:

\begin{equation}
G^{l,t,i}[j] \;=\; \relu{\frac{\partial\, o^{l,t}_{v_t}}{\partial\, A^{l,t}_{i}[-1, j]}},
\qquad j \in \Vset ,
\label{eq:grad}
\end{equation}

\noindent and the same quantity over $j \in \Tset$, written $G^{l,t,i}_{\text{txt}}$.
Each $G^{l,t,i}$ is a map over the $N$ image positions, reshaped to the token
grid $h \times w$. Because
$\partial o / \partial A_{[\cdot,j]} = \langle v_j, g\rangle$, this map carries a
factor of the value-vector norm $\lVert v_j \rVert$ but no factor of $A_j$
itself: DEX-AR reads counterfactual sensitivity, not realised usage.

\section{DEX-AR as published}

\paragraph{Head scores (Eq.~4).} Two scalars per head, each a maximum over the
columns of the respective block. These produce no heatmap; they are a scoring
step.
\begin{align}
S^{l,t,i}_{\text{img}}  &= \max_{j \in \Vset} \; G^{l,t,i}[j], \label{eq:simg}\\
S^{l,t,i}_{\text{txt}} &= \max_{j \in \Tset} \; G^{l,t,i}_{\text{txt}}[j]. \label{eq:stext}
\end{align}

\paragraph{Head filter and per-token map (Eq.~5).} Each head is weighted by how
much more it attends to image than to text, and the weighted maps are summed
over \emph{all} $L$ layers and all $H$ heads:
\begin{align}
w^{l,t,i} &= \relu{S^{l,t,i}_{\text{img}} - S^{l,t,i}_{\text{txt}}}, \label{eq:w}\\
E^{(t)}   &= \sum_{l=1}^{L}\sum_{i=1}^{H} w^{l,t,i} \cdot G^{l,t,i}. \label{eq:etoken}
\end{align}

\paragraph{Token filter and sentence map (Eq.~6).} The per-token maps are
aggregated under a token weight that applies the same difference to the
\emph{global} maxima, taken across every layer and head:
\begin{align}
\delta^{t} &= \relu{\max_{l,i} S^{l,t,i}_{\text{img}} \;-\; \max_{l,i} S^{l,t,i}_{\text{txt}}}, \label{eq:delta}\\
\bar{E}    &= \sum_{t=1}^{T} \delta^{t} \cdot E^{(t)}. \label{eq:esent}
\end{align}

\paragraph{Where the weight comes from.} Both filters are decided by
$\max_p G(p)$ --- a head's single tallest cell, and nothing else about the shape
of its map. Two heads with the same peak score identically whether that peak sits
in a coherent region or alone.

\section{After: the rank prefilter}

\paragraph{The operator.} Let $R_k$ take, at every cell, the $k$-th order
statistic of that cell's $3\times3$ neighbourhood on the token grid, with
replicate padding at the border:
\begin{equation}
\bigl(R_k\,G\bigr)(p) \;=\; \operatorname{sort}\Bigl(\bigl\{\,G(q) \;:\; q \in \mathcal{N}_{3\times3}(p)\,\bigr\}\Bigr)_{k},
\qquad k \in \{0,\dots,8\}.
\label{eq:rank}
\end{equation}
$k = 0$ is erosion, $k = 4$ the median, $k = 8$ dilation. Border cells replicate
rather than zero-pad: zero padding would erode the border as a padding artifact,
which on a model with a border-peak bias would look like the method working.

\paragraph{Why it targets \eqref{eq:simg}.} An isolated spike cannot survive a
sub-median $R_k$ --- its neighbourhood holds eight small values --- while a
coherent region can. Applying $R_k$ \emph{before} \eqref{eq:simg} rewrites the
head score from ``the tallest cell'' into ``the tallest cell that has
neighbours'', without touching the form of Eq.~5 or Eq.~6.

\paragraph{The rescale.} $S_{\text{txt}}$ lives on the one-dimensional token
sequence and receives no matching operation. Left alone, the subtraction in
\eqref{eq:w} would compare a filtered quantity against an unfiltered one and the
head weights would collapse. So $S_{\text{img}}$ is rescaled per $(t,l)$ by the
ratio of medians over heads:
\begin{equation}
c^{l,t} \;=\; \frac{\operatorname{median}_{i}\,\max_p G^{l,t,i}(p)}
                   {\operatorname{median}_{i}\,\max_p \bigl(R_k G^{l,t,i}\bigr)(p)} .
\label{eq:rescale}
\end{equation}

\paragraph{The two lines that change.} Equations \eqref{eq:grad},
\eqref{eq:stext}, \eqref{eq:w}, \eqref{eq:etoken}, \eqref{eq:delta} and
\eqref{eq:esent} are untouched. Only these are new:
\begin{align}
\widetilde{G}^{l,t,i}         &= R_k\, G^{l,t,i}, \tag{\ref{eq:grad}$'$}\\
\widetilde{S}^{l,t,i}_{\text{img}} &= c^{l,t} \cdot \max_{j \in \Vset} \; \widetilde{G}^{l,t,i}[j], \tag{\ref{eq:simg}$'$}
\end{align}
with $\widetilde{G}$ and $\widetilde{S}_{\text{img}}$ substituted for $G$ and
$S_{\text{img}}$ everywhere downstream. Setting $k$ to nothing recovers the
published method exactly: the prefilter is off by default.

\section{Metrics}

All reported $\times 100$. $A$ is the map bilinearly upsampled to mask
resolution and min-max normalised; $M$ is the ground-truth mask (the union of
every object present).

\begin{align*}
\text{soft-IoU} &= \frac{\sum AM}{\sum A + \sum M - \sum AM} &&\text{threshold-free}\\
\text{IoU}      &= \max_{\tau \in \{\tau_1 \dots \tau_{20}\}} \operatorname{IoU}\bigl(A > \tau,\; M\bigr) &&\text{best of $k=20$ thresholds, per prediction}\\
\text{EPG}      &= \frac{\sum AM}{\sum A} &&\text{share of attribution energy inside the mask}\\
\text{SNR}      &= 10 \log_{10}\!\left( \frac{\langle A, M\rangle / \lVert M \rVert_1}{\langle A, 1-M\rangle / \lVert 1-M \rVert_1} \right) &&\text{$0$ for a uniform map}
\end{align*}

\noindent Soft-IoU and EPG both rise with object area, so neither separates
signal from noise on its own: on the 1449-image PascalVOC val split a random
smooth map scores soft-IoU $18.41$ / IoU $25.46$ / EPG $25.24$ / SNR $-0.01$, and
a uniform map scores $25.24$ on all three of the paper's metrics. SNR is the only
one of the four with a defined zero. Recompute the chance lines for any other
split with \texttt{experiments/voc\_nulls.py}.

\section{Where this lives in the code}

\begin{tabular}{@{}ll@{}}
\toprule
\eqref{eq:grad}, \eqref{eq:simg}, \eqref{eq:stext} & \texttt{dexar/wrapper.py}, the per-layer loop in \texttt{compute\_dexar} \\
\eqref{eq:rank}                                    & \texttt{dexar/wrapper.py}, \texttt{DexarWrapper.\_rank\_filter} \\
\eqref{eq:rescale}                                 & \texttt{dexar/wrapper.py}, the \texttt{map\_filter} branch \\
\eqref{eq:w}, \eqref{eq:etoken}                    & \texttt{dexar/wrapper.py}, ``Head filtering (Eq.~5)'' \\
\eqref{eq:delta}, \eqref{eq:esent}                 & \texttt{dexar/wrapper.py}, sentence aggregation \\
metrics                                            & \texttt{experiments/eval\_masks.py}, \texttt{score()} \\
\bottomrule
\end{tabular}

\medskip
\noindent The prefilter is reached as
\texttt{compute\_dexar(..., map\_filter="rank:k")}, or \texttt{MAP\_FILTER=rank:k}
for \texttt{eval\_masks.py}. \texttt{rank:k:noscale} ablates \eqref{eq:rescale}.
\texttt{drop\_delta} sets $\delta^{t} = 1$, which switches off \eqref{eq:delta}.

\end{document}
